\documentclass[11pt]{article}
\usepackage[margin=1in]{geometry}
\usepackage{amsmath,amssymb,amsthm}
\usepackage{array}
\usepackage{parskip}
\usepackage{booktabs}

\emergencystretch=3em

\newtheorem{theorem}{Theorem}
\newtheorem{lemma}[theorem]{Lemma}
\newtheorem{corollary}[theorem]{Corollary}
\newtheorem{proposition}[theorem]{Proposition}
\theoremstyle{definition}
\newtheorem{definition}[theorem]{Definition}
\theoremstyle{remark}
\newtheorem{remark}[theorem]{Remark}

\newcommand{\mex}{\operatorname{mex}}

\title{Disproof of conjecture \texttt{00000001203}}
\author{earthking11}
\date{2026-09-13}

\begin{document}
\maketitle

\section*{Verdict: FALSE}

We disprove conjecture \texttt{00000001203} as stated. The refutation is
unconditional and elementary: for the move set $S = \{1,3\}$ with $k = 3$, the
Sprague--Grundy (SG) sequence is $0,1,0,1,\dots$, whose least period is $2$. Since
$k = 3 \in S$, the conjecture demands period exactly $k+1 = 4$; it is $2$.
Moreover $\max(S) = 3 = k$, so the two possible readings of the parameter $k$
coincide and the refutation is interpretation-free. Corroborating witnesses
($S = \{2\}$, $k = 2$, and the clause-1 witness $S = \{1,2\}$) show that the
failure is not caused by our choice of witness.

\section{The conjecture}

Quoted verbatim from \texttt{conjectures/00000001203.md}:

\begin{quote}
\textbf{Definition:} An addition-subtraction game is a subtraction game whose
legal moves are the fixed values in a set $S$; the SG period is the least
positive period of $g(n)$. \textbf{Conjecture:} If $S \subset \{1,\dots,k\}$,
then the period of $g$ divides $2^k-1$ if and only if $S$ is nonempty and does
not contain $k$; when $k \in S$ the period is exactly $k+1$.
\end{quote}

\noindent (The conjecture is filed bilingually; the Chinese original is quoted
in \texttt{README.md}.) The statement has two clauses:

\begin{enumerate}
\item[(C1)] \emph{(the ``if and only if'')} the least period of $g$ divides
  $2^k - 1$ \quad $\Longleftrightarrow$ \quad $S$ is nonempty and $k \notin S$;
\item[(C2)] \emph{(the special case)} if $k \in S$, the least period is exactly
  $k + 1$.
\end{enumerate}

Both clauses are refuted below. The primary witness fails (C2); the witness
$S=\{1,2\}$ fails (C1) under every reading of $k$.

\section{Subtraction games and the SG recursion}

\begin{definition}[SG sequence]
For a finite set $S \subset \mathbb{N}_{>0}$ of legal moves (an
\emph{addition-subtraction game}), the SG sequence $g \colon \mathbb{N} \to
\mathbb{N}$ is
\[
g(0) = 0, \qquad g(n) = \mex\{\, g(n-s) \;:\; s \in S,\ s \le n \,\} \quad (n \ge 1),
\]
where $\mex(T)$ is the least non-negative integer not in $T$.
\end{definition}

\begin{definition}[SG period]
The \emph{SG period} is the least positive $p$ with $g(n+p) = g(n)$ for all
$n \ge 0$.
\end{definition}

\section{The primary witness: $S = \{1,3\}$, $k = 3$}

\begin{theorem}\label{thm:primary}
Let $S = \{1,3\}$ and $k = 3$. Then the SG sequence is
\[
g(0), g(1), g(2), g(3), g(4), g(5), \dots \;=\; 0, 1, 0, 1, 0, 1, \dots
\]
and $g(n) = n \bmod 2$ for all $n$. Its least period is $2$.
\end{theorem}

\begin{proof}
The recursion reads
\[
g(0)=0,\quad
g(1)=\mex\{g(0)\},\quad
g(2)=\mex\{g(1)\},\quad
g(n)=\mex\{g(n-3),\,g(n-1)\}\ \ (n \ge 3).
\]
We prove $g(n) = n \bmod 2$ by strong induction. The first terms are
\[
g(0)=\mex\varnothing=0,\quad
g(1)=\mex\{0\}=1,\quad
g(2)=\mex\{1\}=0 .
\]
For $n \ge 3$, both predecessors $n-3$ and $n-1$ have the same parity as
$n-1$, so by induction $g(n-3) = g(n-1) = (n-1) \bmod 2 =: a \in \{0,1\}$.
Hence
\[
g(n) = \mex\{a, a\} = \mex\{a\} = 1 - a = 1 - \bigl((n-1) \bmod 2\bigr) = n \bmod 2 .
\]
Thus $g(n+2) = g(n)$ for all $n$, so $2$ is a period, while $g(1) = 1 \ne 0 =
g(0)$ excludes period $1$. Therefore the least period is exactly $2$.
\end{proof}

\begin{corollary}\label{cor:primary}
The primary witness refutes conjecture \texttt{00000001203}. Here
$k = 3 \in S$, so clause (C2) demands least period $k+1 = 4$; the actual least
period is $2 \ne 4$. Equivalently, there is no way for the period to be both
$4$ and the least period, because $2$ already is.
\end{corollary}

\begin{remark}
For this witness clause (C1) happens to be \emph{satisfied}: $k = 3 \in S$ makes
the right-hand side false, and $2 \nmid 2^3 - 1 = 7$ makes the left-hand side
false as well, so the biconditional is true. It is clause (C2) that fails. This
is exactly why we also exhibit a clause-(C1) witness in
Section~\ref{sec:corroborating}.
\end{remark}

\section{The reading of the parameter $k$ is ambiguous --- and the primary
witness is immune}

The conjecture says ``if $S \subset \{1,\dots,k\}$'', which admits two natural
readings of $k$:

\begin{itemize}
\item[(A)] \emph{Ambient bound.} $k$ is the bound in $\{1,\dots,k\}$, so the
  hypothesis is $S \subseteq \{1,\dots,k\}$ for a $k$ that may strictly exceed
  $\max(S)$. This reading is supported by the wording: ``does not contain $k$''
  is a genuine, non-vacuous condition only when $k$ can lie outside $S$.
\item[(B)] \emph{Charitable / tight reading.} $k = \max(S)$.
\end{itemize}

Under reading (B) the condition ``$k \notin S$'' is never satisfiable (because
$\max(S) \in S$ by definition whenever $S \neq \varnothing$), so the ``if and
only if'' degenerates. We therefore treat both readings explicitly.

\begin{proposition}\label{prop:immune}
The primary witness $S = \{1,3\}$, $k=3$ refutes clause (C2) under
\emph{both} readings (A) and (B): under (A) the ambient bound is $k = 3$; under
(B) we have $\max(S) = 3$, so again $k = 3$. Since the two values agree and
$k \in S$, clause (C2) demands period $k+1 = 4$ in either case, while the least
period is $2$.
\end{proposition}

Thus the disproof does not depend on resolving the ambiguity. By contrast, the
earlier candidate $S = \{1\}$ with $k = 3$ is fragile: under reading (B) one has
$k = 1$ and the period $2 = k+1$ satisfies clause (C2), so that witness refutes
the conjecture only under reading (A). We do not rely on it.

\section{Corroborating witnesses}
\label{sec:corroborating}

\subsection*{Second-clause witness: $S = \{2\}$, $k = 2$}

\begin{theorem}\label{thm:secondary}
Let $S = \{2\}$. Then
\[
g = 0, 0, 1, 1, 0, 0, 1, 1, \dots, \qquad g(n) = \Bigl\lfloor \frac{n}{2}
\Bigr\rfloor \bmod 2,
\]
with least period $4$. Since $k = 2 \in S$, clause (C2) demands least period
$k+1 = 3$; but the least period is $4 \ne 3$. Again $\max(S) = 2 = k$, so the
witness is interpretation-free.
\end{theorem}

\begin{proof}
Here $g(n) = \mex\{g(n-2)\}$ for $n \ge 2$, with $g(0)=g(1)=0$ (for $n=1$ the
only move, $2$, is illegal). Since $\mex\{0\}=1$ and $\mex\{1\}=0$, the
sequence alternates $0,1,0,1,\dots$ starting from $g(2)$ but with the initial
double $g(0)=g(1)=0$; hence $g(n) = \lfloor n/2 \rfloor \bmod 2$, which has
least period $4$. (The values are $0,0,1,1$ repeating, so no smaller positive
period works.)
\end{proof}

\subsection*{First-clause witness: $S = \{1,2\}$}

\begin{theorem}\label{thm:clause1}
Let $S = \{1,2\}$. Then $g = 0,1,2,0,1,2,\dots$, i.e.\ $g(n) = n \bmod 3$, and
the least period is $3$. Moreover $3 \mid 2^2 - 1 = 3$.
\end{theorem}

\begin{proof}
$g(0)=0$, $g(1)=\mex\{g(0)\}=1$, and for $n \ge 2$,
$g(n)=\mex\{g(n-2),g(n-1)\}$. The first values are $0,1,2,0,1,2,\dots$: from
$\{0,1,2\}$ the mex of the two preceding distinct residues is the third, and
the pattern repeats cyclically. Hence $g(n) = n \bmod 3$ and the least period is
$3$ (the values $0,1,2$ are pairwise distinct, so periods $1$ and $2$ are
impossible). Finally $2^2 - 1 = 3$ and $3 \mid 3$.
\end{proof}

\begin{corollary}\label{cor:clause1}
Clause (C1) fails for $S = \{1,2\}$ under both readings of $k$.

\emph{Charitable reading (B):} $k = \max(S) = 2$. Since $k \in S$, the
right-hand side \emph{``$S$ is nonempty and $k \notin S$''} is false; clause
(C1) therefore requires the period \emph{not} to divide $2^k - 1 = 3$. But the
period is $3$ and $3 \mid 3$, so the required divisibility holds and the
biconditional fails.

\emph{Ambient reading (A):} take $k = 3$; then $S = \{1,2\}$ is a
\emph{proper} subset of $\{1,2,3\}$, the right-hand side is true ($3 \notin S$),
so clause (C1) requires $3 \mid 2^3 - 1 = 7$; but $3 \nmid 7$, and the
biconditional fails again.
\end{corollary}

\begin{remark}[On the only conceivable escape]
One might try to rescue the conjecture by restricting clause (C2) to the full
sets $\{1,\dots,k\}$ (the only move sets for which $k = \max(S)$ and
$S = \{1,\dots,k\}$ coincide). The text does not state such a restriction, and
it would not help anyway: clause (C1) is refuted by $S = \{1,2\}$
(Corollary~\ref{cor:clause1}), which for ambient $k=3$ is a proper subset and
for the charitable reading has $S \subsetneq \{1,\dots,k\}$ as well.
We also emphasise that the exact reading of $k$ is genuinely ambiguous in the
conjecture text (ambient bound versus $\max(S)$); we have deliberately chosen
the primary witness so that the ambiguity is immaterial.
\end{remark}

\section{Least periods for $k = 2, 3, 4$}

Table~\ref{tab:periods} lists the least period of $g$ for every nonempty
$S \subseteq \{1,\dots,k\}$, together with the conjecture's prediction under
both readings of $k$: reading (A) with ambient bound $k$, reading (B) with
$k = \max(S)$. A ``fails'' entry means that one of the clauses (C1), (C2) is
violated by that row.

\begin{table}[htbp]
\centering
\small
\begin{tabular}{c l c c c}
\toprule
$k$ & $S$ & least period & reading (A): ambient $k$ & reading (B): $k=\max(S)$ \\
\midrule
$2$ & $\{1\}$       & $2$ & fails & holds \\
$2$ & $\{2\}$       & $4$ & fails & fails \\
$2$ & $\{1,2\}$     & $3$ & fails & fails \\
\midrule
$3$ & $\{1\}$       & $2$ & fails & holds \\
$3$ & $\{2\}$       & $4$ & fails & fails \\
$3$ & $\{3\}$       & $6$ & fails & fails \\
$3$ & $\{1,2\}$     & $3$ & fails & fails \\
$3$ & $\{1,3\}$     & $2$ & fails & fails \\
$3$ & $\{2,3\}$     & $5$ & fails & fails \\
$3$ & $\{1,2,3\}$   & $4$ & holds & holds \\
\midrule
$4$ & $\{1\}$       & $2$ & fails & holds \\
$4$ & $\{2\}$       & $4$ & fails & fails \\
$4$ & $\{3\}$       & $6$ & fails & fails \\
$4$ & $\{4\}$       & $8$ & fails & fails \\
$4$ & $\{1,2\}$     & $3$ & holds & fails \\
$4$ & $\{1,3\}$     & $2$ & fails & fails \\
$4$ & $\{1,4\}$     & $5$ & fails & fails \\
$4$ & $\{2,3\}$     & $5$ & holds & fails \\
$4$ & $\{2,4\}$     & $6$ & fails & fails \\
$4$ & $\{3,4\}$     & $7$ & fails & fails \\
$4$ & $\{1,2,3\}$   & $4$ & fails & holds \\
$4$ & $\{1,2,4\}$   & $3$ & fails & fails \\
$4$ & $\{1,3,4\}$   & $7$ & fails & fails \\
$4$ & $\{2,3,4\}$   & $6$ & fails & fails \\
$4$ & $\{1,2,3,4\}$ & $5$ & fails & fails \\
\bottomrule
\end{tabular}
\caption{Least periods of the SG sequence for every nonempty $S \subseteq
\{1,\dots,k\}$, $k = 2,3,4$, and whether conjecture \texttt{00000001203} holds
for that row under each reading of $k$. The primary witness is the row
$k=3$, $S=\{1,3\}$ (period $2$ versus the demanded $2k+1$); the secondary
witness is $k=2$, $S=\{2\}$; the clause-1 witness is $S=\{1,2\}$. Note that
$S=\{1\}$ is the fragile case: it ``holds'' under reading (B) because then
$k=1$ and the period $2 = k+1$.}
\label{tab:periods}
\end{table}

\begin{remark}
Reading (A) is violated in $3$ of $3$ rows for $k=2$, $6$ of $7$ rows for
$k=3$, and $13$ of $15$ rows for $k=4$. Reading (B) is violated in $2, 5, 13$
of those rows respectively. In particular each table contains violating rows
under either reading, so the conjecture is false no matter how $k$ is read.
\end{remark}

\section{Reproducibility}

The finite computations are carried out by \texttt{reproduce.py}, which uses
only the Python~3 standard library:

\begin{quote}\ttfamily
\$ python3 reproduce.py\\
... \textit{(witnesses, all periods, and the violations tables)}\\
PASS: all checks verified; conjecture 00000001203 is FALSE.
\end{quote}

It computes $g$ by the mex recursion for all nonempty $S \subseteq \{1,\dots,k\}$
with $k = 2,3,4$ over $600$ terms and extracts the least period by direct
search. It verifies in particular:

\begin{itemize}
\item $S=\{1,3\}$: $g(n) = n \bmod 2$ for the first $600$ terms, least period
  $2$, whereas the conjecture (as $3 \in S$) demands $k+1 = 4$;
\item $S=\{2\}$: least period $4 \ne 3 = k+1$;
\item $S=\{1,2\}$: least period $3$ with $3 \mid 2^2 - 1$ (clause (C1) fails
  under the charitable reading), and $3 \nmid 2^3 - 1 = 7$ for ambient $k=3$
  with $S$ a proper subset (clause (C1) fails under the ambient reading).
\end{itemize}

It prints \texttt{PASS} and exits $0$ exactly when every check holds, and exits
non-zero otherwise.

\section{Formal verification}

The refutation is also formalised in core Lean~4 (no Mathlib, no \texttt{sorry})
under \texttt{lean4/}. The primary witness is the well-founded recursion

\begin{quote}\ttfamily
def g : Nat $\to$ Nat\\
\hspace*{2em}| 0 => 0\\
\hspace*{2em}| 1 => mex [g 0]\\
\hspace*{2em}| 2 => mex [g 1]\\
\hspace*{2em}| n+3 => mex [g (n+2), g n]
\end{quote}

\noindent with the closed form \texttt{g\_eq : g n = n \% 2} proved by
\texttt{Nat.strongRecOn}, the least-period predicate \texttt{IsLeastPeriod},
and the two key theorems

\begin{quote}\ttfamily
theorem least\_period\_two : IsLeastPeriod 2\\
theorem not\_least\_period\_four : $\lnot$ IsLeastPeriod (3 + 1)
\end{quote}

\begin{sloppypar}
\noindent The clause-1 witness is formalised as \texttt{g12} with
\texttt{least\_period\_three\_12 : IsLeastPeriod12 3} and
\texttt{three\_dvd\_two\_sq\_sub\_one : 3 $\mid$ 2\textasciicircum2 - 1}. All of
these are collected in \texttt{conjecture\_00000001203\_false}. The axiom audit
\texttt{lake env lean Check.lean} reports only \texttt{[propext, Quot.sound]}
and no \texttt{sorryAx}.
\end{sloppypar}

\section{Status against the submission rules}

Rule 3 requires each submission to contain the LaTeX source code, a PDF
document, and a Lean 4 project.

\begin{itemize}
\item \textbf{LaTeX source} --- present (\texttt{main.tex}), a standalone
  \texttt{article} that uses only \texttt{geometry}, \texttt{amsmath},
  \texttt{amssymb}, \texttt{amsthm}, \texttt{array}, \texttt{parskip}, and
  \texttt{booktabs}, and compiles with \texttt{tectonic main.tex}.
\item \textbf{PDF document} --- present at \texttt{build/main.pdf}, produced by
  \texttt{tectonic main.tex} and moved into \texttt{build/}.
\item \textbf{Lean 4 project} --- present under \texttt{lean4/}, pinned to
  \texttt{leanprover/lean4:v4.33.1} with library \texttt{tlmc1203}; builds with
  \texttt{lake build}; the audit \texttt{lake env lean Check.lean} shows no
  \texttt{sorryAx} and no Mathlib dependency.
\item \textbf{Reproduction} --- \texttt{python3 reproduce.py} prints
  \texttt{PASS} and exits $0$.
\end{itemize}

\end{document}
